\documentclass[conference]{IEEEtran}
\usepackage[T1]{fontenc}
\usepackage[utf8]{inputenc}
\usepackage[spanish,es-tabla]{babel}
\usepackage{amsmath}
\usepackage{amssymb}
\usepackage{booktabs}
\usepackage{array}
\usepackage{multirow}
\usepackage{microtype}
\usepackage{tikz}
\usepackage{pgfplots}
\usepackage{hyperref}
\usepackage{enumitem}
\usepackage{siunitx}
\usetikzlibrary{arrows.meta,positioning,shapes.geometric,fit,calc}
\pgfplotsset{compat=1.18}
\hypersetup{hidelinks}

\title{Diseno e Implementacion de Nodos Worker Distribuidos para Procesamiento Paralelo de Imagenes con gRPC y \texttt{asyncio}}

\author{
\IEEEauthorblockN{Autor del proyecto}
\IEEEauthorblockA{Documento tecnico del prototipo IEEE\\
Python, gRPC, \texttt{asyncio}, procesamiento de imagenes}
}

\begin{document}
\maketitle

\begin{abstract}
Este documento describe el diseno y la implementacion de un subsistema de nodos \emph{worker} para procesamiento distribuido de imagenes en paralelo. El prototipo fue desarrollado en Python con \texttt{asyncio} y \texttt{grpc.aio}, y contempla tres nodos worker, un coordinador mock para pruebas locales, cola de prioridad, scheduling interno, control de recursos, backpressure, retries, cancelacion, metricas, healthchecks y persistencia durable de estado y eventos. La propuesta se enfoca exclusivamente en la capa worker y en el plano de comunicacion con el servidor central. El resultado es un prototipo ejecutable, contenedorizado y validado por pruebas automatizadas, orientado a despliegues reales con soporte de almacenamiento compartido, adaptadores de OCR e inferencia y recuperacion ante fallos.
\end{abstract}

\begin{IEEEkeywords}
sistemas distribuidos, gRPC, procesamiento de imagenes, Python, asyncio, scheduling, tolerancia a fallos
\end{IEEEkeywords}

\section{Introduccion}
Los sistemas de procesamiento de imagenes con alta concurrencia requieren desacoplar el plano de coordinacion del plano de ejecucion. En este proyecto se implemento la porcion de \emph{worker nodes} de una arquitectura distribuida donde un servidor central asigna trabajos y los workers ejecutan pipelines de transformacion de imagenes en paralelo. El objetivo no fue construir un backend completo, sino un componente robusto y listo para integrarse con un coordinador externo.

El worker resultante acepta tareas por gRPC, las prioriza internamente, evalua admision segun recursos, ejecuta transformaciones sobre imagenes y reporta estados, progreso, errores y resultados al coordinador. La implementacion incluye soporte para entradas \texttt{payload}, \texttt{input\_path} e \texttt{input\_uri}, asi como salidas locales o sobre \texttt{output\_uri}.

\section{Objetivos y alcance}
El alcance del prototipo cubre:
\begin{itemize}[leftmargin=*]
\item nodos worker distribuidos;
\item contrato gRPC de control y callback;
\item cola de prioridad con repriorizacion;
\item scheduler interno y control de admision;
\item procesamiento paralelo de imagenes;
\item retries, cancelacion y graceful shutdown;
\item metricas, logs y healthchecks;
\item persistencia durable de pendientes, resultados y spool de eventos.
\end{itemize}

Quedaron fuera de alcance: frontend, autenticacion, base de datos completa, backend de negocio y coordinador productivo. Para pruebas locales se implemento un coordinador mock.

\section{Arquitectura general del prototipo}
La Fig.~\ref{fig:architecture} resume la arquitectura ejecutable del prototipo: un coordinador mock centralizado, tres workers con la misma imagen de despliegue y un backend de almacenamiento configurable.

\begin{figure*}[t]
\centering
\begin{tikzpicture}[
  font=\small,
  service/.style={draw, rounded corners=4pt, thick, minimum width=3.3cm, minimum height=1.05cm, align=center, fill=blue!8},
  worker/.style={draw, rounded corners=4pt, thick, minimum width=3.1cm, minimum height=1.15cm, align=center, fill=green!10},
  sub/.style={draw, rounded corners=3pt, thick, minimum width=3.1cm, minimum height=0.7cm, align=center, fill=white},
  store/.style={draw, rounded corners=4pt, thick, minimum width=4.2cm, minimum height=1cm, align=center, fill=orange!10},
  legend/.style={draw, rounded corners=3pt, thick, fill=white, align=left, inner sep=5pt},
  control/.style={-{Latex[length=2.5mm]}, ultra thick, blue!70!black},
  callback/.style={-{Latex[length=2.5mm]}, ultra thick, green!50!black},
  sharedarrow/.style={-{Latex[length=2.3mm]}, thick, gray!70}
]
\node[service] (coord) at (0,4.6) {Coordinador mock\\endpoint gRPC};
\node[worker] (w1) at (-5.2,2.45) {Worker 1\\\texttt{127.0.0.1:50051}};
\node[worker] (w2) at (0,2.45) {Worker 2\\\texttt{127.0.0.1:50061}};
\node[worker] (w3) at (5.2,2.45) {Worker 3\\\texttt{127.0.0.1:50071}};

\node[sub] (w1q) at (-5.2,1.15) {Cola + scheduler + executor};
\node[sub] (w2q) at (0,1.15) {Cola + scheduler + executor};
\node[sub] (w3q) at (5.2,1.15) {Cola + scheduler + executor};

\node[store] (storage) at (-2.6,-0.85) {Storage compartido o local\\filesystem, \texttt{file://}, \texttt{s3://}, MinIO};
\node[store] (obs) at (3.7,-0.85) {Observabilidad\\Prometheus, \texttt{/livez}, \texttt{/readyz}};
\node[legend, anchor=west] (lg) at (-7.05,-2.15) {
\tikz{\draw[control] (0,0) -- (0.9,0);} Plano de control gRPC: \texttt{SubmitTask}, \texttt{GetNodeStatus}, \texttt{CancelTask}, \texttt{DrainNode}, \texttt{ShutdownNode}\\
\tikz{\draw[callback] (0,0) -- (0.9,0);} Callback gRPC: \texttt{ReportProgress}, \texttt{ReportResult}, \texttt{Heartbeat}\\
\tikz{\draw[sharedarrow] (0,0) -- (0.9,0);} Persistencia, resultados, metricas y healthchecks
};

\draw[control] ($(coord.south)+(-0.85,0)$) to[out=-160,in=90] ($(w1.north)+(0.35,0)$);
\draw[control] ($(coord.south)+(0.0,0)$) -- ($(w2.north)+(0.0,0)$);
\draw[control] ($(coord.south)+(0.85,0)$) to[out=-20,in=90] ($(w3.north)+(-0.35,0)$);

\draw[callback] ($(w1.north)+(-0.35,0)$) to[out=90,in=-160] ($(coord.south)+(-1.65,0)$);
\draw[callback] ($(w2.north)+(0.0,0)$) -- ($(coord.south)+(0.0,0)$);
\draw[callback] ($(w3.north)+(0.35,0)$) to[out=90,in=-20] ($(coord.south)+(1.65,0)$);

\draw[sharedarrow] (w1q.south) to[out=-80,in=145] (storage.north west);
\draw[sharedarrow] (w2q.south) -- ($(storage.north)+(0.8,0)$);
\draw[sharedarrow] (w3q.south) to[out=-100,in=35] (storage.north east);

\draw[sharedarrow] (w1.east) to[out=-18,in=160] (obs.west);
\draw[sharedarrow] (w2.east) -- (obs.west);
\draw[sharedarrow] (w3.south) -- ($(obs.north)+(1.05,0)$);
\end{tikzpicture}
\caption{Arquitectura del prototipo. Azul: RPC de control desde el coordinador hacia cada worker. Verde: callbacks del worker hacia el coordinador. Gris: acceso a storage, metricas y healthchecks.}
\label{fig:architecture}
\end{figure*}

Cada worker mantiene los siguientes subsistemas internos:
\begin{itemize}[leftmargin=*]
\item \textbf{servicio gRPC de control}: recibe tareas y comandos del coordinador;
\item \textbf{priority queue}: implementada con \texttt{heapq};
\item \textbf{scheduler local}: combina prioridad, aging, deadline y costo estimado;
\item \textbf{resource manager}: evalua CPU, memoria, I/O y concurrencia activa;
\item \textbf{execution manager}: despacha transformaciones sobre \texttt{ThreadPoolExecutor} y procesos dedicados cuando se requiere aislamiento;
\item \textbf{reporter}: mantiene un canal persistente gRPC hacia el coordinador;
\item \textbf{state store}: preserva resultados terminales, pendientes y eventos pendientes de entrega.
\end{itemize}

\section{Contrato gRPC y flujo de comunicacion}
El contrato de comunicacion fue implementado en \texttt{proto/worker\_node.proto}. El plano de control y el plano de callback quedaron separados como se indica en la Tabla~\ref{tab:rpc}.

\begin{table}[h]
\centering
\caption{RPCs implementados entre coordinador y worker}
\label{tab:rpc}
\begin{tabular}{>{\raggedright\arraybackslash}p{2.7cm} >{\raggedright\arraybackslash}p{4.5cm}}
\toprule
\textbf{RPC} & \textbf{Funcion} \\
\midrule
\texttt{SubmitTask} & Envia una tarea al worker y devuelve aceptacion, duplicado y posicion de cola. \\
\texttt{GetNodeStatus} & Consulta el estado instantaneo del nodo. \\
\texttt{CancelTask} & Cancela una tarea en cola o en ejecucion. \\
\texttt{DrainNode} & Cambia el nodo a modo draining. \\
\texttt{ShutdownNode} & Programa apagado ordenado del worker. \\
\texttt{ReportProgress} & Callback del worker al coordinador con estados intermedios. \\
\texttt{ReportResult} & Callback con el resultado final o el error terminal. \\
\texttt{Heartbeat} & Señal periodica de vida y capacidad del nodo. \\
\bottomrule
\end{tabular}
\end{table}

El ciclo de vida de una tarea se muestra en la Fig.~\ref{fig:lifecycle}. El worker reporta estados \texttt{QUEUED}, \texttt{ADMITTED}, \texttt{RUNNING}, \texttt{RETRY\_SCHEDULED}, \texttt{SUCCEEDED}, \texttt{FAILED}, \texttt{CANCELLED} y \texttt{REJECTED}.

\begin{figure}[h]
\centering
\begin{tikzpicture}[
  node distance=0.55cm and 0.8cm,
  state/.style={draw, rounded corners, minimum width=2.35cm, minimum height=0.7cm, align=center, fill=gray!8},
  arrow/.style={-{Latex[length=2mm]}, thick}
]
\node[state] (queued) {QUEUED};
\node[state, below=of queued] (admitted) {ADMITTED};
\node[state, below=of admitted] (running) {RUNNING};
\node[state, below left=0.75cm and 0.4cm of running] (retry) {RETRY\_SCHEDULED};
\node[state, below=of running] (success) {SUCCEEDED};
\node[state, below right=0.75cm and 0.4cm of running] (failed) {FAILED};
\node[state, right=1.8cm of running] (cancelled) {CANCELLED};
\node[state, left=1.8cm of queued] (rejected) {REJECTED};

\draw[arrow] (queued) -- (admitted);
\draw[arrow] (admitted) -- (running);
\draw[arrow] (running) -- (success);
\draw[arrow] (running) -- (failed);
\draw[arrow] (running) -- (retry);
\draw[arrow] (retry) |- (queued);
\draw[arrow] (queued) -- (rejected);
\draw[arrow] (running) -- (cancelled);
\draw[arrow] (queued) -| (cancelled);
\end{tikzpicture}
\caption{Estados internos reportados durante el ciclo de vida de una tarea.}
\label{fig:lifecycle}
\end{figure}

\section{Cola de prioridad y scheduling interno}
El worker administra una cola de prioridad local sobre \texttt{heapq}. El score implementado combina prioridad estatica, urgencia por deadline, envejecimiento, costo estimado, reintentos y tamano de imagen:
\begin{equation}
\label{eq:score}
\text{score} = w_p P + w_d D + w_a A - w_c C - w_r R - w_s S
\end{equation}
donde $P$ es la prioridad normalizada, $D$ la urgencia relativa al deadline, $A$ el aging acumulado, $C$ el costo estimado, $R$ el numero de intentos previos y $S$ el peso relativo de la imagen.

La politica final es hibrida:
\begin{itemize}[leftmargin=*]
\item la prioridad acelera trabajos criticos;
\item el aging evita starvation;
\item el deadline incrementa el score cuando una tarea se acerca a su vencimiento;
\item el costo estimado favorece una variante de SJF sin sacrificar prioridad.
\end{itemize}

El estimador de costo se actualiza con promedio exponencial:
\begin{equation}
\tau_{n+1} = \alpha t_n + (1-\alpha)\tau_n
\end{equation}
donde $t_n$ es el tiempo real de servicio observado y $\tau_n$ la prediccion anterior.

\section{Control de admision y recursos}
Antes de ejecutar una tarea, el worker evalua si existe capacidad disponible en CPU, memoria, I/O y GPU opcional. La capacidad efectiva se aproxima por:
\begin{equation}
\label{eq:capacity}
C_{\text{ef}} = \min(f_{\text{cpu}}, f_{\text{ram}}, f_{\text{gpu}}, f_{\text{io}})\cdot C_{\text{base}}
\end{equation}

Una tarea solo se despacha cuando sus requerimientos estimados son compatibles con la foto actual de recursos y cuando la cola no ha superado el \emph{high water mark}. En saturacion, el worker deja de aceptar carga nueva y responde \texttt{REJECTED} o mantiene las tareas pendientes en cola.

\section{Modelo de concurrencia en Python}
El diseno de concurrencia se dividio de la siguiente manera:
\begin{itemize}[leftmargin=*]
\item \texttt{asyncio} para networking, control gRPC, scheduler, heartbeats y reporter;
\item \texttt{heapq} para la cola priorizada;
\item \texttt{ThreadPoolExecutor} para las transformaciones integradas del pipeline;
\item procesos dedicados solo para aislamiento explicito o extensiones no cooperativas.
\end{itemize}

El \texttt{ExecutionManager} usa semaforos para limitar la concurrencia activa. Las transformaciones integradas de \texttt{resize}, \texttt{rotate}, \texttt{blur}, \texttt{sharpen} y \texttt{brightness/contrast} fueron implementadas de forma cooperativa, permitiendo cancelacion durante el pipeline sin depender normalmente de matar procesos.

\section{Persistencia durable, errores y recuperacion}
La tolerancia a fallos se resolvio en tres niveles:
\begin{enumerate}[leftmargin=*]
\item \textbf{deduplicacion}: resultados terminales persistidos por \texttt{idempotency\_key};
\item \textbf{restauracion de pendientes}: tareas en cola o programadas para retry se recargan al reinicio;
\item \textbf{spool durable de eventos}: progreso y resultados se almacenan hasta recibir \texttt{ack} del coordinador.
\end{enumerate}

Si el coordinador se desconecta, el worker no pierde informacion: guarda el evento y reintenta la entrega con backoff exponencial y \emph{replay} al restablecer el canal. Si una tarea falla durante la ejecucion, se registra el error y se decide entre \texttt{RETRY\_SCHEDULED} y \texttt{FAILED} segun \texttt{max\_retries}.

El resultado final del worker incluye:
\begin{itemize}[leftmargin=*]
\item estado terminal;
\item ruta o URI de salida;
\item formato de imagen;
\item dimensiones y tamano;
\item metadata adicional del pipeline;
\item codigo y mensaje de error cuando aplica.
\end{itemize}

\section{Backpressure y analisis de capacidad}
El comportamiento esperado ante saturacion puede analizarse con modelos clasicos de colas. Se utilizaron como referencia:
\begin{equation}
L = \lambda W
\end{equation}
\begin{equation}
\rho = \frac{\lambda}{\mu}
\end{equation}
\begin{equation}
W_{M/M/1} = \frac{1}{\mu - \lambda}
\end{equation}

La Fig.~\ref{fig:mm1} ilustra la forma en que crece la espera analitica cuando la utilizacion $\rho$ se aproxima a 1. Aunque la curva es teorica, justifica el uso de backpressure, watermarks y control de admision.

\begin{figure}[h]
\centering
\begin{tikzpicture}
\begin{axis}[
  width=0.95\columnwidth,
  height=5cm,
  xlabel={Utilizacion $\rho$},
  ylabel={Factor de espera relativo},
  ymin=0,
  ymax=12,
  xmin=0.1,
  xmax=0.95,
  domain=0.1:0.92,
  samples=100,
  grid=both
]
\addplot[thick, blue] {1/(1-x)};
\end{axis}
\end{tikzpicture}
\caption{Curva analitica de crecimiento de espera relativa al aproximarse $\rho$ a 1.}
\label{fig:mm1}
\end{figure}

\section{Operaciones soportadas}
Las transformaciones integradas del worker son:
\begin{itemize}[leftmargin=*]
\item escala de grises;
\item \texttt{resize};
\item \texttt{crop};
\item \texttt{rotate};
\item \texttt{flip};
\item \texttt{blur};
\item \texttt{sharpen};
\item brillo/contraste;
\item watermark de texto;
\item conversion de formato JPG, PNG y TIF.
\end{itemize}

Adicionalmente, el worker integra dos extensiones por adaptador:
\begin{itemize}[leftmargin=*]
\item \textbf{OCR}: ejecutable configurable por \texttt{WORKER\_OCR\_COMMAND};
\item \textbf{inferencia}: ejecutable configurable por \texttt{WORKER\_INFERENCE\_COMMAND}.
\end{itemize}

\section{Observabilidad y despliegue}
El prototipo expone:
\begin{itemize}[leftmargin=*]
\item metricas Prometheus;
\item \texttt{/livez} y \texttt{/readyz};
\item logs estructurados;
\item tracing distribuido opcional con OpenTelemetry y exportacion OTLP;
\item mTLS opcional en el servidor gRPC del worker y en el canal hacia el coordinador;
\item despliegue local con \texttt{docker compose};
\item tres instancias worker y un coordinador mock.
\end{itemize}

El estado de readiness considera cola, memoria disponible y conectividad hacia el coordinador. Las metricas mas relevantes incluyen longitud de cola, tiempos de espera, latencia de procesamiento, retry count, utilization de recursos, tareas activas, cancelaciones y conectividad gRPC.

\section{Validacion del prototipo}
El repositorio fue validado con pruebas unitarias e integracion sobre los componentes de:
\begin{itemize}[leftmargin=*]
\item procesamiento de imagenes;
\item persistencia de estado y spool;
\item recuperacion de pendientes;
\item cancelacion;
\item comunicacion worker--coordinador.
\end{itemize}

La ejecucion final de la suite arrojo catorce pruebas aprobadas. Esto confirma que el worker es ejecutable, capaz de aceptar tareas, encolarlas, despacharlas, procesar imagenes, reportar progreso y recuperarse ante fallos comunes del plano de control.

\section{Limitaciones y trabajo futuro}
\noindent\textbf{Limitaciones principales.}
\begin{itemize}[leftmargin=*, itemsep=0.25ex, topsep=0.35ex]
\item la cola sigue siendo local a cada nodo; no se implemento una cola distribuida con propiedad compartida;
\item el coordinador usado en la validacion es un mock y no un backend de negocio completo;
\item OCR e inferencia requieren un backend ejecutable provisto por el operador;
\item el balanceo global entre workers sigue perteneciendo al coordinador y no al worker individual.
\end{itemize}

\noindent\textbf{Trabajo futuro.}
La siguiente etapa consiste en integrar un coordinador productivo, una cola global con ownership distribuido y mecanismos de balanceo entre workers basados en capacidad agregada del cluster.

\section{Conclusiones}
Se implemento exitosamente un subsistema de workers distribuidos para procesamiento de imagenes con Python, \texttt{asyncio} y gRPC. El resultado no es un diagrama teorico sino un prototipo ejecutable con tres nodos, durabilidad local o compartida, control de admision, cancelacion, retries, observabilidad y despliegue en contenedores. La arquitectura lograda separa de forma clara coordinacion, scheduling, ejecucion y reporte, lo que la vuelve apta como base para una evolucion posterior hacia un sistema distribuido productivo.

\end{document}
